\chapter{第十一届大唐杯高职组省赛}
\fancyhead[L]{\color{H6}\kaishu\faIcon{atom}\;第十一届大唐杯高职组省赛}
\fancyhead[R]{\color{H6}\kaishu\rightmark\,}

\date{2024年4月19日}{}{\href{https://github.com/xubenshan/dtcup}{\textbf{才疏学浅的小熊}}}
{}
{https://www.xiaohongshu.com/user/profile/6192219d0000000021028647}{小红书}
{https://space.bilibili.com/692546609}{B站账号}


\section{单选题（共40小题，共120分）}



\begin{choice}{C}[]
    以下不属于AI赋能5G方面应用的是（\quad）
    \begin{tasks}(4)
        \task 流量预测
        \task 异常检测
        \task 机器翻译
        \task 网络优化
    \end{tasks}
\end{choice}


\begin{choice}{A}[]
    不同地物类型的阴影效应的影响大小不一，以下（\quad）地区受阴影效应影响相对最大？
    \begin{tasks}(4)
        \task 密集城区
        \task —般城区
        \task 郊区
        \task 农村
    \end{tasks}
\end{choice}




\begin{choice}{D}[]
    下列（\quad）是虚拟专用网络（VPN）的安全功能？
    \begin{tasks}(4)
        \task 隧道，防火墙和拨号
        \task 验证，访问控制和密码
        \task 压缩，解密和密码
        \task 加密，鉴别和密钥管理
    \end{tasks}
\end{choice}




\begin{choice}{D}[]
    5G NR构架中数据传输的方式不包括
    \begin{tasks}(4)
        \task 回传
        \task 前传
        \task 中传
        \task 后传
    \end{tasks}
\end{choice}


\begin{choice}{D}[]
    5G NR系统中无线承载可分为信令承载和数据承载，以下属于数据承载的为
    \begin{tasks}(4)
        \task SRB0
        \task SRB2
        \task SRB1
        \task DRB
    \end{tasks}
\end{choice}

\begin{choice}{C}[]
    在5G通信系统中，SDN是指
    \begin{tasks}(4)
        \task 网络功能虚拟化
        \task 移动边缘计算
        \task 软件定义网络
        \task 控制功能重构
    \end{tasks}
\end{choice}




\begin{choice}{A}[]
    “云-边-端”智能网联系统中，下列说法正确的是
    \begin{tasks}(1)
        \task OBU主要实现车辆信息的广播，与其他设备的交互
        \task 智能路侧设备实现局部区域交通态势感知
        \task RSU可以实现高精地图服务，提供安全效率类服务
        \task 边缘云主要实现总体交通态势管理，全局调度
    \end{tasks}
\end{choice}


\begin{choice}{B}[]
    AI的英文缩写是
    \begin{tasks}(2)
        \task Automatice Information
        \task Artifical Intelligence
        \task Artifical Information
        \task Automatic Intelligence
    \end{tasks}
\end{choice}


\begin{choice}{B}[]
    一般情况下，EMB6216BBU与AAU电源线规格分别是
    \begin{tasks}(4)
        \task 25方，8方
        \task 16方，10方
        \task 25方，10方
        \task 16方，8方
    \end{tasks}
\end{choice}


\begin{choice}{C}[]
    大唐5G基站设备升级时，基站BBU升级安装包选择下载到基站的文件路径为
    \begin{tasks}(4)
        \task /data
        \task /ata1
        \task /ata2
        \task /cgf
    \end{tasks}
\end{choice}

\begin{choice}{D}[]
    5G基站设备中，当DCPD至AAU电源线长度大于50米，小于等于80米时，直流电源线应使用
    \begin{tasks}(4)
        \task 2×25mm²
        \task 2×10mm²
        \task 2×6mm²
        \task 2×16mm²
    \end{tasks}
\end{choice}


\begin{choice}{B}[]
    5G NR协议栈中，负责QoS流到DRB的映射的协议层
    \begin{tasks}(4)
        \task PDCP
        \task SDAP
        \task RLC
        \task RRC
    \end{tasks}
\end{choice}


\begin{choice}{B}[]
    以下哪个技术不属于大数据技术在5G通信网络中的具体应用
    \begin{tasks}(2)
        \task 大数据收集分析技术
        \task 移动云计算技术
        \task 无线监控技术
        \task 大数据融合技术
    \end{tasks}
\end{choice}


\begin{choice}{B}[]
    5GFR2最大传输带宽配置
    \begin{tasks}(4)
        \task 600MHz
        \task 400MHz
        \task 100MHz
        \task 200MHz
    \end{tasks}
\end{choice}

\begin{choice}{A}[]
    5G网络中所提到的NR指的是
    \begin{tasks}(4)
        \task 5G新空口
        \task 5G演进空口
        \task 4.5G演进空口
        \task 4.5G空口
    \end{tasks}
\end{choice}

\begin{choice}{B}[]
    5G NR帧结构中，一个无线帧由多少个子帧构成
    \begin{tasks}(4)
        \task 5
        \task 10
        \task 2
        \task 8
    \end{tasks}
\end{choice}



\begin{choice}{A}[]
    6G超低卫星网络中，使用的超低轨道卫星是指距离地面高度约（）公里的轨道。
    \begin{tasks}(4)
        \task 300
        \task 400
        \task 450
        \task 350
    \end{tasks}
\end{choice}

\begin{choice}{C}[]
    大唐5G设备开通调测时，采用近端开站的方式，所使用的开站软件是
    \begin{tasks}(4)
        \task ATP
        \task Wireshark
        \task LMT
        \task OMC
    \end{tasks}
\end{choice}


\begin{choice}{D}[]
    在5G网络架构中，以下选项哪一项是AMF的功能

    \begin{tasks}(2)
        \task 会话的建立修改删除
        \task 下行数据的通知
        \task 漫游功能
        \task 注册管理
    \end{tasks}
\end{choice}



\begin{choice}{C}[]
    在5G基站安装时，DCPD与电源柜之间的电源线规格应为

    \begin{tasks}(4)
        \task 10方
        \task 16方
        \task 25方
        \task 8方
    \end{tasks}
\end{choice}



\begin{choice}{B}[]
    5G控制信道资源调度的基本单位CCE与REG之间的正确关系为

    \begin{tasks}(2)
        \task 1CCE=8REG
        \task 1CCE=6REG
        \task 1CCE=9REG
        \task 1CCE=7REG
    \end{tasks}
\end{choice}




\begin{choice}{B}[]
    5G基站设备EMB6216进行19英寸标准机柜安装时，该机柜需提供多少U的安装空间
    \begin{tasks}(4)
        \task 2U
        \task 4U
        \task 6U
        \task 8U
    \end{tasks}
\end{choice}

\begin{choice}{B}[]
    当前5G网络建设中，以下哪个频段分配给中国联通使用
    \begin{tasks}(2)
        \task 2515-2675MHz
        \task 3500-3600MHz
        \task 3400-3500MHz
        \task 2715-2815MHz
    \end{tasks}
\end{choice}

\begin{choice}{B}[]
    大规模天线形成阵列技术的英文简称是
    \begin{tasks}(2)
        \task MIMO
        \task Massive MIMO
        \task MU-MIMO
        \task mMTC
    \end{tasks}
\end{choice}


\begin{choice}{D}[]
    6G上游行业不包括
    \begin{tasks}(4)
        \task 射频器件
        \task 光纤光缆
        \task 芯片
        \task 智能制造
    \end{tasks}
\end{choice}


\begin{choice}{C}[]
    在5G NR网络中终端在初始接入过程中，基站向核心网发的第一条信息是
    \begin{tasks}(2)
        \task Uplink NAS Transfer
        \task RRC Setup Complete
        \task Initial UE Message
        \task UE Capability Information
    \end{tasks}
\end{choice}



\begin{choice}{A}[]
    大唐5G基站安装时，当BBU安装+AAU设备安装2台时，上级输入空开建议为
    \begin{tasks}(2)
        \task 2路80A
        \task 1路63A
        \task 2路100A
        \task 1路80A
    \end{tasks}
\end{choice}


\begin{choice}{A}[]
    下面哪个是工业技术软件化的重要成果，本质上是一种与原宿主解耦的工业技术经验、规律与知识的沉淀、转化和应用的载体。
    \begin{tasks}(4)
        \task 工业APP
        \task 手机APP
        \task 控制系统
        \task 电脑端软件
    \end{tasks}
\end{choice}


\begin{choice}{D}[]
    一个典型的无线对讲通话设备通常不包括
    \begin{tasks}(2)
        \task 麦克风和扬声器
        \task 调频器
        \task 天线系统
        \task 干扰调节器
    \end{tasks}
\end{choice}


\begin{choice}{B}[]
    5G帧结构中采用常规CP时，一个时隙中固定有多少个符号
    \begin{tasks}(4)
        \task 12
        \task 14
        \task 20
        \task 6
    \end{tasks}
\end{choice}


\begin{choice}{A}[]
    我国发布的V2X标准定义了5大类消息来满足不同场景的需求，其中不包括
    \begin{tasks}(4)
        \task RRC
        \task SPAT
        \task BSM
        \task MAP
    \end{tasks}
\end{choice}


\begin{choice}{A}[]
    5G基站BBU设备EMB6216机框，仅有2块基带板板卡的情况下，应当优先插几槽位？
    \begin{tasks}(2)
        \task 2和3槽位
        \task 2和7槽位
        \task 2和5槽位
        \task 2和6槽位
    \end{tasks}
\end{choice}



\begin{choice}{A}[]
    在5G 移动通信每十年一代，6G技术其实是5G代际更新的一个新技术，预期其标准化制定时间
    \begin{tasks}(4)
        \task 2025年
        \task 2029年
        \task 2030年
        \task 2024年
    \end{tasks}
\end{choice}


\begin{choice}{C}[]
    5G核心网架构类型为
    \begin{tasks}(4)
        \task 星型
        \task 矩阵型
        \task 总线型
        \task 网状型
    \end{tasks}
\end{choice}


\begin{choice}{B}[]
    SA组网场景下，站点开通时5G基站连接到核心网的SCTP链路，该链路的协议类型为

    \begin{tasks}(4)
        \task XNAP
        \task NGAP
        \task X2
        \task AMF
    \end{tasks}
\end{choice}


\begin{choice}{B}[]
    5G NR空中接口协议层中，Uu口控制面最高层是
    \begin{tasks}(4)
        \task MAC
        \task RRC
        \task PHY
        \task PDCP
    \end{tasks}
\end{choice}


\begin{choice}{C}[]
    DC供电系统的电压波动范围为
    \begin{tasks}(4)
        \task -40V至-53V
        \task -40V至-50V
        \task -40V至-57V
        \task -40V至-55V
    \end{tasks}
\end{choice}


\begin{choice}{C}[]
    5G NR系统中，5G基站和核心网通过哪个接口相连
    \begin{tasks}(4)
        \task X2
        \task Xn
        \task NG
        \task S1
    \end{tasks}
\end{choice}


\begin{choice}{B}[]
    5G NR中，SSS在时域上占用SSB的第几个符号位
    \begin{tasks}(4)
        \task 4
        \task 3
        \task 2
        \task 1
    \end{tasks}
\end{choice}


\begin{choice}{B}[]
    5G基站安装时，为避免自激，GPS放大器安装位置应距离GPS天线至少多少米
    \begin{tasks}(4)
        \task 10
        \task 20
        \task 30
        \task 40
    \end{tasks}
\end{choice}



\section{多选题（共30小题，共120分）}

\begin{choice}{\;ABC\;}[]
    5G R15版本，PDCCH信道CCE聚合等级可能为
    \begin{tasks}(4)
        \task 4
        \task 1
        \task 2
        \task 3
    \end{tasks}
\end{choice}

\begin{choice}{\;ACD\;}[]
    移动通信系统中，覆盖问题包括弱覆盖和
    \begin{tasks}(2)
        \task 导频污染
        \task 天线覆盖
        \task 覆盖空洞
        \task 越区覆盖
    \end{tasks}
\end{choice}

\begin{choice}{\;AB\;}[]
    根据项目的不同阶段，勘察可分为哪两个阶段
    \begin{tasks}(2)
        \task 复勘
        \task 初勘
        \task 天面勘察
        \task 网络环境勘察
    \end{tasks}
\end{choice}

\begin{choice}{\;BCD\;}[]
    在5G网络中基站间通过XN接口相连，则属于XN接口协议栈组成的为
    \begin{tasks}(4)
        \task SDAP
        \task GTP-U
        \task SCTP
        \task XN-AP
    \end{tasks}
\end{choice}


\begin{choice}{\;ABD\;}[]
    以下哪些技术属于大数据技术在5G通信网络中的具体应用

    \begin{tasks}(2)
        \task 大数据融合技术
        \task 大数据收集分析技术
        \task 移动云计算技术
        \task 无线监控技术
    \end{tasks}
\end{choice}

\begin{choice}{\;ABCD\;}[]
    属于5G超密集组网应用场景有

    \begin{tasks}(2)
        \task 大型活动场馆
        \task 商业中心
        \task 密集住宅区
        \task 交通枢纽
    \end{tasks}
\end{choice}

\begin{choice}{\;ABD\;}[]
    关于GPS的安装规范，以下哪些说法是正确的

    \begin{tasks}(1)
        \task 放大器安装在离GPS天线50m-150m处，安装在墙上或走线架上
        \task 一定要在避雷针的45度防雷保护范围内
        \task GPS安装不需要接地
        \task GPS天线宜远离其他发射或接收设备，不要安装在微波天线、高压线下方，避免发射天线的辐射方向对准GPS天线。
    \end{tasks}
\end{choice}


\begin{choice}{\;ABCD\;}[]
    关于HSCTA板卡功能，下列说法正确的是

    \begin{tasks}(1)
        \task 基站系统与GPS的同步功能
        \task 提供与AAU连接的Ir接口
        \task L2处理功能
        \task 与BBU内部各板卡之间的业务、信令交换处理功能
    \end{tasks}
\end{choice}


\begin{choice}{\;ABD\;}[]
    移动通信系统中，常见干扰类型包括
    \begin{tasks}(2)
        \task 阻塞干扰
        \task 发射机交调干扰
        \task 绕射干扰
        \task 杂散干扰
    \end{tasks}
\end{choice}


\begin{choice}{\;BCD\;}[]
    基站设备安装时，GPS的安装环境要求有
    \begin{tasks}(1)
        \task GPS天线位于避雷针保护范围内，最好是区域内的最高点
        \task 一般来说GPS天线安装建议选取位置较开阔，天空可视性较好，垂直方向无阻挡且便于安装的位置
        \task 严禁将GPS天线安装在基站等系统的辐射天线主瓣面内，不能和全向天线安装在同水平面内
        \task 尽量将GPS天线安装在安装地点的南边，周围对天线的遮挡不超过30度，天线竖直向上的视角应大于120度
    \end{tasks}
\end{choice}

\begin{choice}{\;AD\;}[]
    在基站勘察设计过程中，以下哪些不属于勘察流程
    \begin{tasks}(2)
        \task 仔细研读设备说明书
        \task 天馈勘察
        \task 勘察前准备
        \task 基站设备安装
    \end{tasks}
\end{choice}

\begin{choice}{\;BCD\;}[]
    6G网络将在智享生活、智赋生成等方面催生全新的应用场景，这些场景将需要以下哪些全新的网络功能和能力
    \begin{tasks}(2)
        \task 毫秒级的时延体验
        \task 太比特级的峰值速率
        \task 超过1000km/h的移动速度
        \task 数字李生
    \end{tasks}
\end{choice}

\begin{choice}{\;BCD\;}[]
    5G NR系统中，MIB消息中包含以下哪些内容

    \begin{tasks}(2)
        \task PHICH配置
        \task SIB1的PDCCHCORESET配置
        \task 系统帧号
        \task 半帧指示
    \end{tasks}
\end{choice}

\begin{choice}{\;ABC\;}[]
    在工业互联网体系架构中，网络层技术发挥着至关重要的作用，以下关于工业互联网网络层技术描述正确的选项是
    \begin{tasks}(1)
        \task 工业以太网、现场总线技术以及无线传感器网络都属于工业互联网网络层的技术范畴
        \task 工业互联网网络层技术需要与IT（信息技术）和OT（操作技术）领域进行深度融合以满足工业自动化和智能化的需求
        \task 工业互联网网络层负责实现设备间的互联互通，包括有线和无线通信技术
        \task 在工业互联网中，网络层技术不关注数据传输的安全性和实时性
    \end{tasks}
\end{choice}

\begin{choice}{\;ACD\;}[]
    5G NR中，RRC状态包括哪些
    \begin{tasks}(2)
        \task RRC CONNECTED
        \task URA CONNECTED
        \task  RRC IDLE
        \task RRC INACTIVE
    \end{tasks}
\end{choice}

\begin{choice}{\;ABD\;}[]
    5G网络中物理下行参考信号DM-RS可能伴随那些物理信道传制
    \begin{tasks}(2)
        \task PDSCH
        \task PBCH
        \task PUCCH
        \task PDCCH
    \end{tasks}
\end{choice}

\begin{choice}{\;AD\;}[]
    决策树主要解决以下哪些问题

    \begin{tasks}(2)
        \task 分类
        \task 以上都不是
        \task 聚类
        \task 回归
    \end{tasks}
\end{choice}

\begin{choice}{\;ABCD\;}[]
    基站前期勘察的准备事项包括
    \begin{tasks}(2)
        \task 工具准备
        \task 人员准备
        \task 信息准备
        \task 拟定计划
    \end{tasks}
\end{choice}

\begin{choice}{\;ABCD\;}[]
    NR系统中，关于SDN技术的特点，说法正确的是
    \begin{tasks}(2)
        \task 控制器软件可编程
        \task 控制平面集中化
        \task 控制转发分离
        \task 转发平面通用化
    \end{tasks}
\end{choice}

\begin{choice}{\;ABD\;}[]
    在5G网络中，UPF集成了核心网的所有用户面功能，UPF通过N3接口连接基站，以下属于N3接口协议栈的为
    \begin{tasks}(2)
        \task IP
        \task UDP
        \task SDAP
        \task GTP-U
    \end{tasks}
\end{choice}

\begin{choice}{\;ABCD\;}[]
    人工智能全方位赋能6G网络安全，6G网络安全内生具备哪些特征
    \begin{tasks}(2)
        \task 虚拟共生
        \task 主动免疫
        \task 泛在协同
        \task 弹性自治
    \end{tasks}
\end{choice}

\begin{choice}{\;ABD\;}[]
    一个典型的无线对讲通话设备通常包括

    \begin{tasks}(2)
        \task 调频器
        \task 天线系统
        \task 干扰调节器
        \task 麦克风和扬声器
    \end{tasks}
\end{choice}

\begin{choice}{\;ACD\;}[]
    移动通信系统中一般快衰落可以细分为
    \begin{tasks}(2)
        \task 空间选择性衰落
        \task 多径选择性衰落
        \task 时间选择性衰落
        \task 频率选择性衰落
    \end{tasks}
\end{choice}


\begin{choice}{\;BC\;}[]
    目前无线网络性能的测试主要工作不包括

    \begin{tasks}(2)
        \task 网络调整测试
        \task 基站开通
        \task 勘察设计
        \task 拉网摸底测试
    \end{tasks}
\end{choice}


\begin{choice}{\;AD\;}[]
    产品设计反馈优化场景中，工业互联网平台可以将哪两个数据反馈到设计和制造阶段，从而改进设计方案，加速创新迭代
    \begin{tasks}(2)
        \task 产品运行数据
        \task 物流数据
        \task 制造数据
        \task 用户使用行为数据
    \end{tasks}
\end{choice}


\begin{choice}{\;ACD\;}[]
    大唐5G基站开通后，关于板卡状态，说法正确的是
    \begin{tasks}(2)
        \task 管理状态应为解锁定状态
        \task 板卡过程状态为未初始化状态
        \task 板卡过程状态应为初始化状态
        \task 运行状态应为正常
    \end{tasks}
\end{choice}


\begin{choice}{\;ABCD\;}[]
    人工智能的基本技术包括
    \begin{tasks}(4)
        \task 机器学习
        \task 知识表示
        \task 推理技术
        \task 搜索技术
    \end{tasks}
\end{choice}


\begin{choice}{\;ABCD\;}[]
    5G NR中无线承载可分为信令承载和数据承载，其中信令承载包括哪些
    \begin{tasks}(4)
        \task SRB3
        \task SRB1
        \task SRB0
        \task SRB2
    \end{tasks}
\end{choice}


\begin{choice}{\;ABCD\;}[]
    6G网络秉承兼容、跨域、分布、内生、至简、李生六大设计理志进行架构设计，以下说法正确的是
    \begin{tasks}(1)
        \task 6G架构设计将由外挂式向内生式转变
        \task 6G架构设计将由集中规划式向分布自治式转变，满足大规模组网下的海量连接和极致性能要求
        \task 6G架构设计将支持固定、移动、卫星等多种接入
        \task 移动通信网络沿着IP化、云化、服务化的方向发展变革
    \end{tasks}
\end{choice}


\begin{choice}{\;BCD\;}[]
    以下选项关于5G帧结构描述正确的是
    \begin{tasks}(2)
        \task 每个subframe的slot个数固定
        \task 特殊子帧位置灵活配置
        \task slot及符号长度与scs有关
        \task 每个subframe的slot个数与scs相关
    \end{tasks}
\end{choice}


\section{判断题（共30小题，共60分）}

\begin{choice}{\;正确\;}[]
    Boosting方法是将几个弱分类器结合成一个强学习器的集成方法。
\end{choice}


\begin{choice}{\;错误\;}[]
    5G基站设备EMB6216的高度是3U。
\end{choice}

\begin{choice}{\;正确\;}[]
    LMT接口是BBU和本地维护终端连接的接口。
\end{choice}

\begin{choice}{\;正确\;}[]
    5G系统中通过小区搜索的目的为与基站取得下行同步。
\end{choice}



\begin{choice}{\;错误\;}[]
    一个典型的无线对讲通话设备通常包括干扰调节器。

\end{choice}


\begin{choice}{\;正确\;}[]
    无线通信可以传送电报、电话、传真、图像、数据及广播和电视节目等通信业务。

\end{choice}

\begin{choice}{\;错误\;}[]
    在5G NR中，NZC代表随机接入前导序列长度，取值可能为838和138。

\end{choice}

\begin{choice}{\;正确\;}[]

    按照5G物理网络架构，前传是用来连接AAU和BBU设备。
\end{choice}


\begin{choice}{\;正确\;}[]
    以技术为主，以研发部为中心的传统产品开发模式，是面向技术先进性而进行的产品开发，不是以客户为中心的产品开发模式。

\end{choice}

\begin{choice}{\;正确\;}[]
    5G NR系统中，初始随机接入失败后，终端可以提升发射功率，继续发起随机接入过程。

\end{choice}


\begin{choice}{\;正确\;}[]
    NR系统中：n41频段的SSB子载波间隔可以为15KHz或者30KHz。

\end{choice}


\begin{choice}{\;正确\;}[]
    不要用力拉扯光纤，或用脚及其它重物踩压光纤，以免造成光纤的损坏。

\end{choice}


\begin{choice}{\;正确\;}[]
    浪涌保护器也要做室内接地。

\end{choice}


\begin{choice}{\;错误\;}[]
    通信按照复用方式的不同可以分为频分、时分和码分三种。

\end{choice}


\begin{choice}{\;正确\;}[]
    移动通信系统中，信令传递通过使用交织，可增强其可靠性。

\end{choice}

\begin{choice}{\;正确\;}[]
    5G不同组网方式类别中option2为独立组网方式。

\end{choice}


\begin{choice}{\;正确\;}[]
    C-6G将会是一个整合了地面无线和卫星的全连接世界。

\end{choice}


\begin{choice}{\;正确\;}[]
    在5G基站设备EMB6116中，HSCTD板卡可以放在0槽位或者1槽位上。

\end{choice}


\begin{choice}{\;错误\;}[]
    5G NR帧结构Normal CP一个子帧固定包含2个时隙。

\end{choice}


\begin{choice}{\;正确\;}[]
    5G基站设备EMB6116能够支持的时钟同步方式包括北斗、GPS、1588v2等。

\end{choice}


\begin{choice}{\;正确\;}[]
    5GSA场景下，业务转发的配置是SMF的功能。

\end{choice}


\begin{choice}{\;错误\;}[]
    华为Mate60Pro的卫星通话需要将信号先发给卫星系统的地面站进行通信。

\end{choice}


\begin{choice}{\;正确\;}[]
    5GR15版本小区带宽sub-6GHz包括：5/10/15/20/25/30/40/50/60/80/100MHz。

\end{choice}

\begin{choice}{\;正确\;}[]
    移动通信系统中，研究传播模型可以通过无线传播理论出发进行研究，也可以通过大量测试数据的基础上统计分析出传播损耗的数学规律进行研究。

\end{choice}

\begin{choice}{\;正确\;}[]
    HBPOD板卡能够支持时钟同步和分发功能。

\end{choice}

\begin{choice}{\;错误\;}[]
    在5G的无线部署中，要求CU和DU必须分离。

\end{choice}

\begin{choice}{\;错误\;}[]
    噪声系数可以衡量接收机、放大器等射频器件的射频性能，它是指输出端信噪比与输入端信噪比之比。

\end{choice}

\begin{choice}{\;正确\;}[]
    6G将具备的感知功能是指利用AI、数字李生等技术进行建模分析和自动决策，提升网络运营效率和系统性能。

\end{choice}

\begin{choice}{\;正确\;}[]
    车联网环境下，两车交互数据只能通过PC5接口交换。

\end{choice}

\begin{choice}{\;错误\;}[]
    5G标准可以分为两个阶段，R17版本属于第二阶段。

\end{choice}







